\clearpage
\section{АНАЛИЗ ПРЕДМЕТНОЙ ОБЛАСТИ И СУЩЕСТВУЮЩИХ АНАЛОГОВ}

\subsection{Особенности биржевых котировок как временных рядов}

Биржевые котировки относятся к классу финансовых временных рядов, для которых характерно сочетание трендовых участков, шумовой компоненты, изменения волатильности и зависимости от внешних новостных и макроэкономических факторов. В отличие от многих инженерных временных рядов, в рыночных данных редко выполняется предположение о строгой стационарности. Распределение доходностей может меняться во времени, амплитуда колебаний возрастает в периоды стрессов, а закономерности, найденные на одном интервале, часто деградируют на другом. Это явление в прикладных работах связывают с изменением рыночного режима и эффектом \textit{concept drift} [3, 25, 26].

С практической точки зрения биржевой ряд представляет собой не просто последовательность значений цены, а комплексную динамическую систему, где существенную роль играют доходности, скорость изменения цены, локальный тренд, объемы торгов и взаимодействие нескольких временных масштабов. Для целей построения торговых стратегий важны не только ожидаемый уровень цены, но и вероятность благоприятного движения в ближайшем будущем, вероятность ложного сигнала, длительность удержания позиции и ожидаемая величина просадки. Вследствие этого выбор модели должен учитывать как прогностическую силу признаков, так и последующее влияние предсказаний на торговый контур [3, 6].

Дополнительная сложность заключается в том, что на финансовых рынках высокая ошибка в несколько отдельных моментов способна перечеркнуть полезность большого числа корректных прогнозов. Например, стратегия может правильно предсказывать большую долю небольших движений, но потерять значительную часть капитала на редких сильных разворотах. Поэтому интерпретация качества модели через одну лишь среднюю ошибку или точность классификации является неполной. Для корректного исследования необходимо переходить от анализа ряда как статистического объекта к анализу результата как динамики стоимости портфеля [7, 27, 28].

\subsection{Классические методы прогнозирования финансовых временных рядов}

Классические методы прогнозирования финансовых временных рядов включают наивные стратегии, модели сглаживания, авторегрессионные модели и сезонные модификации, такие как ARIMA и SARIMA. Их достоинство состоит в хорошо разработанном математическом аппарате, интерпретируемости параметров и понятной процедуре оценки. Для многих экономических и производственных рядов такие модели остаются надежным инструментом. Однако в отношении биржевых котировок их применимость ограничивается тем, что динамика цен часто содержит выраженную нелинейность, структурные сдвиги и слабую устойчивость локальных закономерностей [2, 4].

Наивная модель \textit{random walk}, предполагающая, что лучшее предсказание следующего значения равно текущему, долгое время использовалась как базовая точка отсчета для оценки более сложных алгоритмов. В рамках торговли ей соответствует \textit{Buy \& Hold}, если актив находится в долгосрочном восходящем тренде, или пассивное отсутствие позиции, если цель исследования сводится к прогнозированию самого уровня цены. Скользящие средние, экспоненциальное сглаживание и ARIMA-подходы позволяют выделять локальные зависимости, но плохо переносят резкие изменения режима, характерные для финансовых рынков [4, 25, 26].

В контексте настоящей работы классические статистические модели важны по двум причинам. Во-первых, они задают нижнюю и среднюю планку для сравнения более современных методов. Во-вторых, они демонстрируют ограничение постановки, в которой цель исследования сводится исключительно к прогнозу следующего значения ряда. Даже хорошо подогнанная ARIMA-модель не определяет напрямую, когда открывать позицию, как долго ее удерживать и как учитывать цену ошибки в терминах реального капитала [2, 4, 6].

\subsection{Методы машинного обучения для прогнозирования котировок}

В отличие от классических статистических моделей, методы машинного обучения ориентированы на работу с признаковым представлением задачи. Входом для модели становится не сам ряд в исходном виде, а набор признаков, включающий лаги, доходности, скользящие средние, индикаторы импульса, оценки волатильности и другие характеристики. Такой переход позволяет формулировать задачу как бинарную классификацию направления следующего движения или как регрессию доходности [7, 15].

Для задач табличного прогнозирования биржевых котировок особенно распространены \textit{Logistic Regression}, \textit{Random Forest}, \textit{Gradient Boosting}, а также продвинутые реализации бустинга, такие как CatBoost и XGBoost. Логистическая регрессия обеспечивает высокую интерпретируемость и позволяет быстро получить базовый ориентир качества. Ансамблевые методы, в свою очередь, способны захватывать нелинейные зависимости между признаками и часто демонстрируют более устойчивый результат на шумных данных [13, 14].

Однако переход к машинному обучению не устраняет фундаментальную проблему исследования: прогноз направления сам по себе не гарантирует построение эффективной стратегии. Одна и та же модель может давать приемлемый ROC-AUC, но формировать слишком частые сигналы, приводящие к избыточным комиссиям, либо слабо работать в периоды рыночного боковика. Поэтому в работе ML-модели рассматриваются не как конечный инструмент принятия решений, а как промежуточный источник торгового сигнала, который затем преобразуется в стратегию и оценивается уже по финансовым метрикам [6, 7].

\subsection{Нейросетевые модели для анализа временных рядов}

Нейросетевые модели получили широкое распространение в задачах анализа временных рядов благодаря способности автоматически извлекать сложные нелинейные закономерности из последовательностей данных. Для финансовых задач наиболее часто используются LSTM и GRU, а также их модификации, комбинируемые с attention-механизмами, сверточными блоками или трансформерными архитектурами. Преимущество таких моделей состоит в том, что они потенциально могут учитывать более длинный контекст, чем классические табличные модели [11, 12].

Тем не менее в прикладных задачах алгоритмической торговли нейросетевые подходы нередко сталкиваются с рядом проблем. Во-первых, они требуют более аккуратной настройки гиперпараметров и, как правило, большего объема данных. Во-вторых, их устойчивость на нестационарных финансовых рядах не гарантирована: модель может хорошо описывать один временной режим и деградировать после его смены. В-третьих, интерпретируемость результатов оказывается заметно ниже, чем у линейных и ансамблевых моделей [7, 11, 12].

В настоящей работе LSTM и GRU рассматриваются как важные аналоги в теоретическом обзоре, поскольку они занимают промежуточное положение между классическим ML и RL-подходами. Они по-прежнему ориентированы на прогноз будущего значения или класса, но делают это в более гибкой нелинейной форме. Однако основной экспериментальный контур работы сфокусирован на baseline-стратегиях, трех ML-моделях и RL-агентах, что обусловлено необходимостью сохранить воспроизводимость и умеренную вычислительную сложность на локальной вычислительной платформе [11, 12].

\subsection{Обучение с подкреплением в задачах алгоритмической торговли}

Обучение с подкреплением принципиально отличается от классического прогностического подхода тем, что объектом оптимизации является не ошибка прогноза, а последовательность действий агента в среде. В задачах алгоритмической торговли это особенно важно: вместо предсказания отдельного значения цены агент получает состояние рынка, выбирает действие и получает награду, связанную с изменением стоимости портфеля, комиссией и риском. Таким образом, модель обучается непосредственно в терминах поведения, а не только в терминах прогноза [1, 6, 8, 9].

Наиболее известные RL-алгоритмы, применяемые в торговых задачах, можно разделить на value-based и policy-based семейства. DQN относится к первой группе и аппроксимирует функцию ценности действий, что удобно для дискретных действий \textit{buy}, \textit{sell}, \textit{hold}. PPO относится ко второй группе и обучает параметризованную политику с ограничением шага обновления, обеспечивая более устойчивую оптимизацию. A2C сочетает идеи актор-критик и служит компактной альтернативой при меньших вычислительных затратах [8, 9, 10].

Несмотря на привлекательность RL, его использование в финансовых задачах не гарантирует автоматического превосходства над более простыми подходами. Качество RL-агента крайне зависит от дизайна среды, набора признаков, функции вознаграждения и выбранного рыночного режима. Именно поэтому RL в настоящем исследовании оценивается не как заведомо лучший метод, а как один из классов алгоритмов, потенциально полезный для оптимизации последовательных решений, но требующий строгого экспериментального сопоставления с альтернативами [1, 8, 9].

\subsection{Обзор аналогов и существующих решений: FinRL, TensorTrade, gym-anytrading}

Среди готовых исследовательских и прикладных решений для RL-трейдинга можно выделить платформы FinRL, TensorTrade и библиотеку gym-anytrading. FinRL предоставляет набор стандартных сред, примеров и сценариев обучения RL-агентов на финансовых данных. Его достоинство состоит в ориентации на воспроизводимость и наличии типовых пайплайнов для получения котировок, подготовки данных и запуска алгоритмов. Однако при этом исследователь во многом зависит от архитектурных решений, заложенных в сам фреймворк [21].

TensorTrade ориентирован на более модульный подход. В нем торговая среда строится из взаимодействующих компонентов: источника данных, механизма исполнения, портфеля, действия и функции награды. Такой подход обеспечивает гибкость, но одновременно повышает порог входа и усложняет тонкую настройку исследования. Для учебной и дипломной работы это может привести к значительным трудозатратам на инфраструктурный уровень вместо концентрации на собственно методологическом сравнении [22].

Библиотека gym-anytrading занимает противоположную позицию. Она предлагает минималистичные среды, совместимые с Gym-интерфейсом, что удобно для быстрого прототипирования. Ее слабой стороной является ограниченная выразительность модели рынка и меньшее число готовых сценариев для серьезного экспериментального расширения. По этой причине в настоящей работе сделан выбор в пользу собственной среды на базе Gymnasium: она сохраняет совместимость со stable-baselines3, но при этом остается достаточно прозрачной и контролируемой для дипломного исследования [20, 23].

\subsection{Сравнительная таблица аналогов}

Для систематизации рассмотренных подходов в таблице~\ref{tab:analogs_comparison} приведено сравнительное описание стратегий, моделей и фреймворков, имеющих наибольшее значение для поставленной задачи. Таблица показывает, что ни один из подходов не является универсальным. Простые baseline-стратегии обеспечивают максимальную интерпретируемость, статистические модели подходят для формального анализа временных рядов, ML-модели хорошо работают на табличных признаках, а RL-алгоритмы наиболее естественно описывают последовательное принятие торговых решений.

\begin{center}
\small
\begin{longtable}{|p{2.6cm}|p{2.3cm}|p{3.6cm}|p{3.8cm}|p{3.2cm}|}
\caption{Сравнение существующих аналогов и подходов}\label{tab:analogs_comparison}\\
\hline
\textbf{Подход / средство} & \textbf{Класс} & \textbf{Сильные стороны} & \textbf{Ограничения} & \textbf{Роль в исследовании} \\
\hline
\endfirsthead
\hline
\textbf{Подход / средство} & \textbf{Класс} & \textbf{Сильные стороны} & \textbf{Ограничения} & \textbf{Роль в исследовании} \\
\hline
\endhead
Buy \& Hold & baseline-стратегия & Простота, прозрачность, отражение рыночного бета-эффекта & Не адаптируется к рыночным режимам, не управляет риском активно & Базовый эталон абсолютной доходности \\
\hline
Moving Average Crossover & техническая стратегия & Интерпретируемость, учет тренда, низкая вычислительная сложность & Запаздывание сигнала, чувствительность к параметрам окна, ложные входы на боковом рынке & Сопоставление RL с простым правилом на индикаторах \\
\hline
ARIMA / SARIMA & статистическое прогнозирование & Формальный аппарат, интерпретируемые параметры, удобство для линейной динамики & Плохо описывает нелинейность и смену режимов, требует стационаризации & Аналог прогнозирования ценового ряда \\
\hline
Logistic Regression & классическое МО & Простота, устойчивость, возможность анализа вклада признаков & Ограниченная нелинейность, зависимость от качества признаков & Базовая ML-модель для бинарного прогноза направления \\
\hline
Random Forest & ансамблевое МО & Работа с нелинейными зависимостями, устойчивость к шуму, малая потребность в масштабировании & Снижение интерпретируемости, риск переобучения при слабой валидации & Сравнение с линейной моделью и RL \\
\hline
Gradient Boosting & ансамблевое МО & Высокая гибкость, качественная аппроксимация сложных закономерностей & Требовательность к настройке, чувствительность к качеству данных & Сильный ML-базис в экспериментальном стенде \\
\hline
CatBoost / XGBoost & градиентный бустинг & Высокое качество на табличных данных, эффективная работа с признаками & Выше вычислительные затраты, необходимость тщательного тюнинга & Аналог продвинутого табличного бустинга \\
\hline
LSTM / GRU & нейросетевые модели & Учет последовательной структуры, моделирование временной зависимости & Чувствительность к объему данных, сложность настройки, риск нестабильного обучения & Аналоги глубокого обучения для временных рядов \\
\hline
DQN & RL-алгоритм value-based & Обучение дискретным действиям, прямое связывание состояния и торгового решения & Нестабильность при плохой reward function, чувствительность к коррелированным данным & RL-аналог с дискретным пространством действий \\
\hline
PPO & RL-алгоритм policy-based & Устойчивое обновление политики, хорошая практическая применимость, удобство для continuous optimization & Более высокая стоимость обучения, зависимость от дизайна среды & Основной RL-алгоритм для сравнения риск-скорректированной эффективности \\
\hline
A2C & actor-critic & Более компактная альтернатива PPO, обучает политику и value function совместно & Меньшая устойчивость по сравнению с PPO на шумных рынках & Рассматривается как альтернативный RL-подход в обзоре \\
\hline
FinRL & фреймворк & Готовые среды, примеры RL-трейдинга, акцент на reproducibility & Сильная зависимость от встроенных сценариев, сложность адаптации под конкретный pipeline & Аналог исследовательской платформы \\
\hline
TensorTrade & фреймворк & Модульная архитектура, абстракции обмена, портфеля и вознаграждения & Повышенный порог входа, необходимость глубокой настройки & Аналог конструкторного инструментария \\
\hline
gym-anytrading & библиотека сред & Простота, совместимость с Gym, быстрое прототипирование & Ограниченная гибкость и сравнительно простая рыночная логика & Аналог минималистичной RL-среды \\
\hline
\end{longtable}
\end{center}


Дополнительно в таблице~\ref{tab:methods_overview} показано, что принципиальное различие между классами методов состоит в типе целевого результата и в критерии качества. Для baseline- и RL-подходов естественной метрикой является динамика капитала, тогда как для статистических, ML- и нейросетевых моделей часто сначала оценивается ошибка прогноза, а уже затем финансовая результативность стратегии, построенной на основе этого прогноза.

\begin{table}[H]
\caption{Сопоставление классов методов, входов и критериев оценки}
\label{tab:methods_overview}
\begin{center}
\small
\begin{tabularx}{\textwidth}{|p{2.8cm}|p{2.7cm}|X|p{2.8cm}|}
\hline
\textbf{Группа методов} & \textbf{Входные данные} & \textbf{Целевой результат} & \textbf{Основной критерий оценки} \\
\hline
Baseline-стратегии & Цена закрытия, скользящие средние & Последовательность торговых сигналов по фиксированному правилу & Финансовые метрики стратегии \\
\hline
Статистические модели & Ряд цен или доходностей & Прогноз следующего значения или доходности & Ошибка прогноза и прибыльность стратегии на основе сигнала \\
\hline
Классические ML-модели & Табличные признаки на основе OHLCV и индикаторов & Класс направления движения или доходность следующего периода & ML-метрики и финансовые метрики после конвертации прогноза в стратегию \\
\hline
Нейросетевые модели & Последовательности наблюдений, лаги, индикаторы & Прогноз направления, уровня цены или доходности & Ошибка прогноза, устойчивость и финансовый результат \\
\hline
RL-агенты & Окно признаков, позиция, доля cash & Действие buy / sell / hold & Накопленная награда и финансовые метрики стратегии \\
\hline
\end{tabularx}
\end{center}

\end{table}

\subsection{Постановка проблемы и выводы по главе}

Проведенный анализ предметной области и аналогов позволяет сформулировать центральную проблему исследования более точно. Существующие работы и инструменты часто сосредоточены либо на точности прогноза, либо на отдельных аспектах торговой результативности. Однако для практической задачи алгоритмической торговли этого недостаточно. Высокая точность классификации не тождественна прибыльности стратегии, а высокая доходность на одном интервале не гарантирует устойчивости на другом. Следовательно, требуется единый экспериментальный стенд, который позволяет сравнивать baseline-стратегии, модели машинного обучения и RL-агентов не изолированно, а в одном методическом контуре.

В рамках такой постановки работа решает именно проблему корректного сравнительного анализа. Она не стремится доказать, что RL всегда превосходит классические подходы. Напротив, исследовательская гипотеза формулируется осторожно: RL может быть полезен тогда, когда задача требует оптимизировать последовательность действий и балансировать между доходностью и риском, однако его реальная ценность должна подтверждаться сравнением с простыми и интерпретируемыми альтернативами. Такой вывод задает логический переход к следующей главе, в которой описываются проектирование метода и построение экспериментального стенда для проведения указанного сравнения.
